% TOP_PROBLEM: {"id": 6200062, "title": "Boundaries of Groups and Kleinian Groups — Problem 62", "category_id": 6, "category": {"id": 6, "name": "geometry", "display_name": "Geometry", "description": "Euclidean and non-Euclidean geometry, geometric structures.", "slug": "geometry", "order_index": 6, "created_at": "2026-07-31T15:26:25.671Z"}, "status": "open", "problem_number": "AMR-061-0062", "set_id": 13, "statusQualification": "This defines the standalone conformal-dimension problem and its related attainment question. The source also refers to the preceding group-action conjecture (UnsolvedMath 6200061), concerning an infimum of visual dimensions over geometric actions; that separate conjecture is not identified here with attainment in the full quasisymmetric gauge."}
% FIRST_POSTED: 2026-10-02
% ENTRY_KIND: statement_only
% UnsolvedMath ID 6200062; source catalogue code AMR-061-0062
% !TeX program = lualatex
% Definition and sources only; this file is not an LLM solution attempt.
\documentclass[11pt,a4paper]{article}
\usepackage[margin=25mm]{geometry}
\usepackage{fontspec}
\setmainfont{lmroman10-regular.otf}[BoldFont=lmroman10-bold.otf,
 ItalicFont=lmroman10-italic.otf,BoldItalicFont=lmroman10-bolditalic.otf]
\usepackage{amsmath,amssymb,mathtools}
\usepackage{unicode-math}
\setmathfont{latinmodern-math.otf}
\AtBeginDocument{\let\setminus\smallsetminus}
\usepackage{xurl,hyperref}
\hypersetup{colorlinks=true,urlcolor=blue}
\setcounter{secnumdepth}{0}
\setlength{\parindent}{0pt}
\setlength{\parskip}{5pt}
\setlength{\emergencystretch}{3em}
\begin{document}
\section*{Boundaries of Groups and Kleinian Groups — Problem 62}\label{6200062}
{\small UnsolvedMath ID 6200062 · AMR-061-0062 · Geometry · AMR Open Problem Lists}
\subsection{Definitions and mathematical statement}
Let $\mathcal S\subset[0,1]^2$ be the standard Sierpi\'nski carpet: subdivide a square into nine equal squares, remove the open central square, and repeat this operation in every retained square. Give the intersection of the retained stages its Euclidean distance $d$.

A metric $d'$ on $\mathcal S$ belongs to its quasisymmetric gauge if the identity $(\mathcal S,d)\to(\mathcal S,d')$ is a homeomorphism whose distance ratios are controlled by a fixed distortion homeomorphism. The metric $d'$ is Ahlfors $Q$-regular if it has a Borel measure $\mu$ and a constant $C\geq1$ satisfying
\[
 C^{-1}r^Q\leq\mu(B_{d'}(x,r))\leq Cr^Q
 \quad(0<r\leq\operatorname{diam}_{d'}\mathcal S).
\]
Determine the Ahlfors-regular conformal dimension
\[
 \operatorname{Cdim}_{AR}(\mathcal S)=
 \inf\{Q:\text{an Ahlfors $Q$-regular metric lies in this gauge}\}.
\]
Also determine whether this infimum is attained by such a metric. The Euclidean Hausdorff dimension $\log8/\log3$ is an upper bound, not an asserted value of the infimum.

The metric gauge is essential: replacing the carpet by any homeomorphic compactum with an arbitrary metric changes the optimization problem. Exact determination requires a matching lower bound and construction or limiting argument. Attainment is a separate issue from knowing the numerical infimum; a sequence of admissible metrics with dimensions approaching that number need not supply a metric realizing it.
\subsection{Short English statement}
What is the Ahlfors-regular conformal dimension of the standard carpet, and does a metric attain it?
\subsection{Sources}
\begin{itemize}
\item[S1] ULAM.AI and the UnsolvedMath Contributors, supplied problems.json, record 6200062 (AMR-061-0062), AMR Open Problem Lists. Catalogue identity and source transcription; CC BY 4.0.
\url{https://huggingface.co/datasets/ulamai/UnsolvedMath}.
\item[S2] Kapovich - Problems on Boundaries of Groups and Kleinian Groups (2005), Problem 62, PDF page 17. Original source indexed by AMR-061-0062.
\url{https://www.math.ucdavis.edu/~kapovich/EPR/problems.pdf}.
\item[S3] Anttila and Eriksson-Bique, Cartesian products of Sierpiński carpets do not attain their conformal dimension (2026). The single-carpet attainment question remains distinct from the product result.
\url{https://arxiv.org/abs/2604.19600}.
\end{itemize}
\end{document}
